\chapter*{Abstract}
\addcontentsline{toc}{chapter}{Abstract}

\noindent
This thesis aims to reconstruct snow water equivalent (SWE) across the Italian Hydrological portion of the Greater Alpine Region 
(IH-GAR) using high-resolution temperature and precipitation data. To this end, a degree-day model was developed to simulate both 
snow accumulation and melt processes. In parallel, an observational snow depth (HS) dataset was assembled through a data 
rescue activity and its main numerical issues were subsequently addressed, to ensure the reliability of these HS series. This 
step proved essential for the optimization of model parameters, which was based on comparisons between simulated SWE and observed 
snow depth time series. The calibrated model was then used to generate daily SWE maps for the 1951--2023 period, thereby allowing 
for robust trend estimates. The results show that climate change is significanty affecting both the amount and persistence of snow 
cover. In particular, the annual peak SWE volume in the IH-GAR is decreasing by \SI{0.567}{\giga\meter\cubed} per decade, corresponding to 
approximately \num{5.6} \% of the 1952--1981 climatological mean. The decline is even more pronounced towards the end of the snow 
season, with the volume of water stored as snow on 1 April decresing at a rate of \num{6.5} \% per decade. These trends are 
accompanied by a reduction of snow cover duration, mainly driven by earlier snowmelt during the spring months. As global warming 
continues, further loss are expected. A possible extension of this work would be to investigate future changes in SWE under 
different climate change scenarios. 

\bigskip
\noindent\textbf{Keywords:} Snow Water Equivalent (SWE), degree-day model, climate change, snow depth.

%\begin{otherlanguage}{italian}

%\end{otherlanguage}
